% ============================================================
%  第 1 章  绪论
% ============================================================
\section{绪论}
\label{sec:intro}

\subsection{研究背景与选题意义}

Web 应用已成为承载社会关键业务的基础设施，其安全性直接关系到数据资产与业务连续性。
开放式全球应用安全项目（OWASP）在其发布的 2021 年度十大风险清单中，将\emph{注入}
（Injection）列为第三位风险，而在此前的 2013 与 2017 两个版本中，注入类漏洞长期
位居首位\cite{owasp2021top10}。以 SQL 注入为代表的注入攻击之所以历经二十余年仍未
消退，其根源在于一个结构性缺陷：应用程序将\emph{不可信的用户输入}与\emph{可信的
指令结构}混合于同一文本通道中，致使解释器无法区分二者边界。

针对这一威胁，工业界形成了两条主要防线。其一是\emph{开发侧的根本性修复}，即通过
参数化查询将指令结构与数据彻底分离；其二是\emph{运行侧的旁路防护}，即部署 Web
应用防火墙（Web Application Firewall, WAF）对请求流量实施检测与阻断。前者虽属治本
之策，但在存量系统改造成本高昂、第三方组件不可控、开发规范执行不彻底等现实约束下，
难以做到全覆盖；后者因而成为不可或缺的补充。

然而，当前主流 WAF 的技术路线存在一个值得深究的方法论问题。以 ModSecurity 核心
规则集（Core Rule Set, CRS）为代表的产品，其判定依据是对请求载荷施加正则表达式
匹配\cite{owasp2023crs}。该范式的判定对象是载荷的\textbf{文本形态}，而非其
\textbf{语义内涵}。二者的错配构成了防护体系的根本性弱点：攻击者可通过 URL 多重
编码、HTML 实体编码、注释拆分、数据库方言内联注释、大小写混淆等手段，在\emph{完全
不改变载荷语义}的前提下任意改变其文本形态。防护方因而被迫陷入“新绕过手法出现—
补充规则—再被绕过”的被动循环，规则集规模持续膨胀而防护有效性并未同步提升，
误报率反而随规则数量增长而恶化。

本课题正是针对上述方法论困境而设立。其核心追问是：\emph{能否找到一种检测依据，
使其对攻击者可自由操纵的维度（文本形态）不敏感，而对攻击者必须改变的维度
（语义结构）高度敏感？}

\subsection{国内外研究现状}

\subsubsection{基于规则匹配的检测方法}

规则匹配是工业界应用最广的方案。ModSecurity CRS 通过数千条正则规则覆盖已知攻击
特征，具备部署简便、可解释性强的优点\cite{owasp2023crs}。但如前所述，该方案的
判定对象与攻击本质错配，且规则的通用性与精确性存在难以调和的矛盾：规则过宽则
误报率高，过窄则易被变形绕过。文献\cite{cui2018survey}对该类方法的系统性评估表明，
其对未知变形载荷的检出率显著低于对已知特征的检出率。

\subsubsection{基于词法分析的检测方法}

Galbreath 提出的 libinjection 是该方向的代表性工作\cite{galbreath2012libinjection}。
该方法将载荷切分为词法单元（token）序列，通过有限状态机识别注入特有的 token 组合
模式。相较正则匹配，词法分析对大小写与空白等表层混淆具备天然免疫力，且检测速度
极快。其局限在于\emph{仅停留于 token 序列层面而未构建完整语法结构}，因而难以表达
“条件被剥离”“子查询嵌套”等需要层次信息才能刻画的语义变化，也无法定位攻击发生的
具体结构位置。

\subsubsection{基于语法树验证的检测方法}

Buehrer 等人在 2005 年提出的 SQLGuard 是本课题最直接的理论渊源
\cite{buehrer2005sqlguard}。该工作的核心洞察是：合法查询的语法树结构在运行时应
保持稳定，注入则会引入结构变化；因此可在运行时比对“仅含开发者意图的查询”与
“混入用户输入后的查询”两者的语法树，若不一致则判定为注入。Boyd 与 Keromytis 提出的
SQLrand 采用了另一种思路，通过对 SQL 关键字施加随机化编码，使攻击者注入的标准
关键字无法被解析器识别\cite{boyd2004sqlrand}。Halfond 与 Orso 提出的 AMNESIA 则
结合静态分析与运行时监控，先由静态分析推导各查询点的合法 SQL 模型，再在运行时
校验实际语句是否符合该模型\cite{halfond2006amnesia}。Su 与 Wassermann 从形式语言
理论出发，将注入攻击形式化定义为“用户输入跨越了语法树的非终结符边界”，为该类方法
奠定了理论基础\cite{su2006essence}。

上述工作确立了“语义优于文本”的正确方向，但普遍存在三方面不足：其一，多数方案
要求侵入式改造应用源码（如 SQLGuard 需开发者显式标注查询模板），部署代价高；
其二，判定结果为二值的“匹配/不匹配”，\emph{无法定位注入发生的具体语法节点}，
对安全取证的支撑不足；其三，缺少面向工程落地的完整系统设计，如基线的自动学习与
审核、检测结果的可视化呈现、以及与现有部署形态的兼容。

\subsubsection{基于机器学习的检测方法}

近年来，卷积神经网络\cite{luo2019deep}与人工神经网络\cite{tang2020detection}等
方法被引入注入检测领域，在特定数据集上取得了较高的准确率。此类方法的优势在于
可自动学习特征而无需人工构造规则，但也面临三重制约：模型判定过程缺乏可解释性，
难以向安全运营人员说明“为何判定为攻击”；训练数据的分布偏差会直接转化为检测盲区；
以及推理开销通常高于确定性算法，在高吞吐场景下形成瓶颈。文献\cite{li2020injection}
对该方向的综述亦指出，可解释性不足是其在生产环境落地的主要障碍。

\subsubsection{运行时应用自我保护技术}

运行时应用自我保护（Runtime Application Self-Protection, RASP）通过将检测逻辑
植入应用进程内部，使防护能够观测到应用的真实运行时上下文\cite{wang2022rasp}。
相较部署于流量入口的 WAF，RASP 的关键优势在于\emph{可见性}：它能够观测到已完成
拼接的真实 SQL，而非仅仅是 HTTP 参数。这一优势对判定的确定性具有决定性意义，
构成了本文双层架构设计的直接依据。

\subsubsection{现有方案的对比}

\cref{tab:related-work} 从五个维度对上述方案进行了横向比较。可以看出，现有工作
尚未形成一个同时具备“非侵入部署、确定性判定、结构级定位、可解释输出”四项特性的
完整方案，这正是本课题的切入点。

\begin{table}[htbp]
\centering
\caption{主流 Web 攻击检测方案的横向比较}
\label{tab:related-work}
\small
\begin{threeparttable}
\begin{tabular}{@{}l c c c c c@{}}
\toprule
\textbf{方案类别} & \textbf{判定对象} & \makecell{\textbf{抗编码}\\\textbf{混淆}} &
\makecell{\textbf{结构级}\\\textbf{定位}} & \makecell{\textbf{可}\\\textbf{解释性}} &
\makecell{\textbf{部署}\\\textbf{侵入性}} \\
\midrule
正则规则集\cite{owasp2023crs}        & 文本形态   & 弱   & 无 & 强 & 无 \\
词法状态机\cite{galbreath2012libinjection} & Token 序列 & 中   & 无 & 强 & 无 \\
语法树验证\cite{buehrer2005sqlguard} & 语法结构   & 强   & 无 & 强 & 高 \\
关键字随机化\cite{boyd2004sqlrand}   & 关键字编码 & 强   & 无 & 中 & 高 \\
静态分析建模\cite{halfond2006amnesia}& 语法模型   & 强   & 无 & 中 & 中 \\
神经网络\cite{tang2020detection}     & 统计特征   & 中   & 无 & 弱 & 无 \\
\midrule
\textbf{本文方法}                    & \textbf{语法结构} & \textbf{强} &
\textbf{有} & \textbf{强} & \textbf{无}\tnote{a} \\
\bottomrule
\end{tabular}
\begin{tablenotes}[flushleft]\footnotesize
\item[a] 本文通过字节码注入实现探针植入，无需修改应用源码，故部署侵入性为“无”。
\end{tablenotes}
\end{threeparttable}
\end{table}

\subsection{本文主要工作与创新点}

本文围绕“以语义结构而非文本形态作为检测依据”这一核心思想，完成了从理论方法到
工程系统的完整设计与实现。主要工作与创新点如下。

\paragraph{（1）提出 SQL 抽象语法树结构指纹方法。}
设计了“词法分析—语法分析—字面量归一化—规范化序列化—哈希”的处理流水线，将任意
SQL 语句映射为仅反映其语义结构的定长指纹。本文形式化证明了该指纹同时满足参数
无关性、形态无关性与结构敏感性三条性质（详见\cref{sec:algorithms}），因而对编码
混淆类绕过手法具有\emph{原理上的}免疫能力，而非依赖对已知手法的枚举。

\paragraph{（2）提出基于路径签名多重集差分的结构定位算法。}
针对现有语法树方案“只能判定是否变化、无法定位变化位置”的不足，本文摒弃复杂度过高
的树编辑距离算法\cite{zhang1989simple}，转而采用路径签名多重集差分，以 $O(n)$ 的
时间复杂度精确定位注入所在的语法树节点，并将差分结果映射为具有安全语义的攻击手法
分类（恒真式绕过、联合查询注入、堆叠查询等），为取证与可视化提供确定性依据。

\paragraph{（3）设计网关与 RASP 探针双层联动架构。}
明确区分两层的\emph{可见性}差异：网关层仅能观测 HTTP 参数文本，判定必然是启发式的；
探针层可观测已完成拼接的真实 SQL，判定是确定性的。两层通过统一追踪标识关联，
可完整还原攻击证据链。该架构在保持非侵入部署的同时，兼顾了拦截时效性与判定准确性。

\paragraph{（4）实现完整的工程系统并进行严格验证。}
系统包含检测引擎、RASP 探针、安全网关、控制台服务、受控靶场与可视化前端六个模块，
代码规模约 1.7 万行。采用\emph{三轮测试法}进行验证：分别确认漏洞真实存在、防护
检出有效、代码修复有效。设计 19 项测试用例覆盖 OWASP Top 10 的七个类别，全部
高危漏洞获得防护，平均检测时延 0.93~ms。

\subsection{论文组织结构}

本文余下部分组织如下。\cref{sec:requirements} 进行需求分析，采用 STRIDE 方法
建立威胁模型并以 DREAD 模型完成风险定级。\cref{sec:design} 阐述系统总体设计，
重点论证双层架构的必要性。\cref{sec:algorithms} 是本文的技术核心，给出结构指纹
与结构差分算法的形式化定义、性质证明与复杂度分析。\cref{sec:security} 论述系统
自身的安全设计，包括纵深防御体系、防篡改审计与安全编码实践。
\cref{sec:implementation} 介绍关键模块的实现细节。\cref{sec:evaluation} 报告
实验设计与评估结果，包含功能验证、绕过对抗实验、误报测试与性能测试。
\cref{sec:management} 记录项目管理过程。\cref{sec:conclusion} 总结全文，
并诚实讨论方法的局限性与后续改进方向。
